% !TeX program = xelatex
% Build with: latexmk -xelatex paper_template.tex
\documentclass[11pt,a4paper]{article}
\usepackage[margin=2.5cm]{geometry}

% The first build creates a .bib file. Replace the placeholder in that file.
\begin{filecontents*}{\jobname.bib}
@article{mbgn_placeholder,
  author  = {Weiqi Xue},
  title   = {{MBGN} baseline paper title},
  journal = {Journal name pending},
  year    = {YYYY}
}
\end{filecontents*}

\title{Multimodal Defect Event Recognition in GMAW Guided by Thermal History and Spatiotemporal Registration}
\author{Weiqi Xue\\Department and Institution}
\date{}

\begin{document}
\maketitle

\begin{abstract}
[State the challenge of detecting undercut events in gas metal arc welding (GMAW).]
[Describe how the method aligns welding current and voltage with weld pool images and accounts for thermal history and uncertainty in event boundaries.]
[Report event recall and false alarms per metre of evaluated weld on held-out welds.]
[Summarize the main finding and the conditions under which it holds.]
\end{abstract}
\noindent\textbf{Keywords:} gas metal arc welding, undercut, thermal history, spatiotemporal registration, multimodal fusion

\section{Introduction}
Introduce the MBGN baseline~\cite{mbgn_placeholder}. State which aspects of undercut detection may depend on the duration of the electrical signal window or on its alignment with the images.

\section{Data and Event Annotation}
\subsection{Welding Conditions and Data Acquisition}
Describe the welds, process settings, electrical sampling rates, camera frame rate, and synchronization procedure.
\subsection{Undercut Intervals and Evaluated Weld Length}
Define the normal and undercut classes. Explain how event boundaries were annotated and which weld intervals were eligible for evaluation.

\section{Methods}
\subsection{Multiscale Thermal History}
Define instantaneous electrical power and the features used to estimate heat accumulation over each preceding time window.
\subsection{Causal Spatiotemporal Registration}
Explain how image timestamps map to weld positions, which electrical samples are available at prediction time, and how uncertainty in the annotated boundaries is handled.
\subsection{Multimodal Fusion and Event Detection}
Describe how the model combines electrical signals and images and converts consecutive predictions into undercut events. Figure~\ref{fig:overview} summarizes this procedure.

\begin{figure}[htbp]
  \centering
  \fbox{\parbox[c][32mm][c]{0.85\linewidth}{\centering Insert the method overview here.}}
  \caption{Alignment and fusion of welding current, voltage, and weld pool images for undercut event detection.}
  \label{fig:overview}
\end{figure}

\section{Experimental Design}
\subsection{Weld-Level Splits and Evaluation Metrics}
Identify the training, validation, and test welds. Specify how the detection threshold is selected, how predicted events are matched to annotations, and how evaluated weld length is measured.
\subsection{Baseline and Ablation Studies}
Describe the MBGN baseline and the comparisons used to isolate the effects of thermal history, registration, and boundary uncertainty.

\section{Results and Discussion}
Report the number of annotated, detected, missed, and falsely predicted events on held-out welds. Compare event recall and false alarms per metre across methods and examine representative errors.

\section{Conclusions}
State which research questions the results resolve and which welding conditions the conclusions cover.

\bibliographystyle{unsrt}
\bibliography{\jobname}
\end{document}
